\subsubsection{Ansatz Design}\label{sec:ansatz-design}

An Ansatz is not just any parameterized quantum circuit. It must be carefully designed so that:
\begin{enumerate}
    \item The \hl{solution can be represented}. The first requirement is \textbf{expressibility}. A highly expressive ansatz can generate many different quantum states.
    \item The \hl{circuit can run on real hardware}. The second requirement is \textbf{hardware compatibility}. The ansatz should use gates that the hardware can implement efficiently. This is crucial because every additional gate can introduce noise, errors, and decoherence issues, which is exactly the problem of NISQ devices.
    \item The \hl{optimizer can find the solution}. The third requirement is \textbf{shallow circuit depth}. A shallow circuit is less prone to noise and can be optimized more easily by the classical optimizer. Deep circuits may increase the chances of noise, gate errors, and decoherence.
    
    \textcolor{DarkOrange3}{\faIcon{balance-scale} \textbf{The Fundamental Tradeoff.}} There is a tension between \emph{expressibility} and \emph{noise}. If the ansatz is too simple, the search space may not contain the optimal solution. However, if the ansatz is too complex, it may be too noisy to run on NISQ devices.
    
    Therefore, designing an ansatz is a delicate balance between these two factors. The \hl{goal is to find an ansatz that is expressive enough to capture the solution while being shallow enough to be executed on real hardware without excessive noise.}
\end{enumerate}
If any of these fail, the VQA may fail.

\highspace
\begin{flushleft}
    \textcolor{DarkOrange3}{\faIcon{balance-scale} \textbf{Problem-Agnostic vs Problem-Specific}}
\end{flushleft}
There are two main approaches to ansatz design:
\begin{itemize}
    \item \definition{Problem-Agnostic Ansatz}. This approach uses a \textbf{general-purpose ansatz that is not tailored to any specific problem}. Examples include the hardware-efficient ansatz, which consists of layers of parameterized single-qubit rotations and entangling gates. The \textbf{advantage} of this approach is that it can be \textbf{applied to a wide range of problems without modification}. However, it may \textbf{require more parameters} and \textbf{deeper circuits} to achieve good performance, which can be problematic on NISQ devices.
    \item \definition{Problem-Specific Ansatz}. This approach designs an \textbf{ansatz that is tailored to the specific problem being solved}. For example, in quantum chemistry, one can use the Unitary Coupled Cluster (UCC) ansatz, which is based on classical coupled cluster theory. The \textbf{advantage} of this approach is that it can be \textbf{more efficient} and require \textbf{fewer parameters} than a problem-agnostic ansatz. However, it may require \textbf{more domain knowledge} and may \textbf{not be applicable to other problems}.
\end{itemize}

\highspace
In summary, Ansatz design is the process of choosing a parameterized quantum circuit suitable for a variational algorithm. An ideal ansatz should be shallow, compatible with the available hardware, and sufficiently expressive to represent the solution of the problem. Since the ansatz defines the search space explored by the optimizer, it must be capable of generating the desired solution for some choice of parameters. However, greater expressibility usually comes at the cost of deeper circuits and increased noise, creating a tradeoff between accuracy and hardware limitations.